\chapter{\Large \texorpdfstring{$HM\_QCv$}{HM\_QCv}: A Computationally Efficient Hybrid Model for Enhancing Quality Clusters in Visualization}



Traditional techniques \cite{tsne} \cite{umap_reference} \cite{sun2023laplacian} exist for visualizing high-dimensional data. Still, they face limitations by covering all characteristics, such as (i) minimizing computational costs, (ii) providing high-quality cluster visualizations, and (iii) preserving both local and global structures. Local structures refer to the relationships between nearby data points \cite{tsne}, whereas global structures encompass the overall distribution of the data \cite{umap_reference} \cite{tenenbaum2000global}.

Sun et al. \cite{sun2023laplacian} developed an approach that uses the t-SNE \cite{tsne} methodology to preserve the local and global structure and incorporates a symmetric Laplacian matrix, the regularization technique for cluster quality when visualizing the data. However, this approach, while improving the representation of clusters, also increases the computational complexity due to t-SNE \cite{tsne}.

The existing methodology LAPtSNE \cite{sun2023laplacian}, relies on the t-SNE \cite{tsne} method to compute the distance between the data objects, which is computationally intensive.  To address this limitation, we adopted UMAP \cite{umap_reference} in our proposed approach.  Before adopting UMAP \cite{umap_reference}, we conducted experiments on three benchmark datasets to evaluate both the t-SNE \cite{tsne} and UMAP \cite{umap_reference} approaches, measuring their processing times. The results clearly demonstrate that UMAP \cite{umap_reference} is more computationally efficient than t-SNE \cite{tsne}. Hence, our approach, $HM\_QCv$, is more computationally efficient than LAPtSNE \cite{sun2023laplacian} while ensuring both local and global structure by preserving high-quality cluster visualization, offering promising practical implications for the field of data visualization.


The $HM\_QCv$ approach is designed to preserve both (i) cluster information and (ii) local and global structures of high-dimensional data in visualization. Our approach leverages the strengths of UMAP \cite{umap_reference} and LAPtSNE \cite{sun2023laplacian} while integrating their partial components to develop more comprehensive solutions. This combined method may reveal insights that would be unattainable with either technique used independently. By leveraging UMAP \cite{umap_reference} and LaptSNE \cite{sun2023laplacian}, our approach ensures computational efficiency and significantly reduces processing time compared to LAPtSNE \cite{sun2023laplacian}.